\section{Scope of the Project}
\label{sec:intro_scope}

Based on the literature review so far, this project bounds itself by one governing rule: work proceeds only with attack surfaces for which a dedicated, reusable, oracle-backed benchmark already exists in the reviewed literature, rather than constructing new benchmarks independently. This decision keeps the evaluation grounded and reproducible. The scope may narrow further once early implementation work provides more information about what is achievable within the remaining timeline. Table~\ref{tab:in-scope-surfaces} summarises the attack-surface families currently considered in scope and the benchmarks that govern any claim made about them.

\begin{table}[htbp]
	\centering
	\small
	\caption{Attack surfaces currently considered in scope, and their candidate governing benchmarks (preliminary).}
	\label{tab:in-scope-surfaces}
	\begin{tabular}{@{}p{3.1cm}p{5.8cm}p{4.8cm}@{}}
		\toprule
		\textbf{Attack surface}       & \textbf{Candidate benchmarks}                                                                                                                                                                        & \textbf{Representative vulnerability types}                                                 \\
		\midrule
		Web application (HTTP/HTML)   & Fang et al.\ 15-CVE sandbox \citep{fang2024oneday}; HPTSA 14-CVE zero-day suite \citep{zhu2024hptsa}; CVE-Bench (40 CVEs) \citep{zhu2025cvebench}; PentestEval (346 tasks) \citep{yang2025pentesteval} & SQLi, XSS, CSRF, SSRF, SSTI, LFI, file-upload RCE, IDOR, authentication bypass, JWT forgery \\
		GraphQL APIs                  & PrediQL 6-API suite \citep{liu2025prediql}                                                                                                                                                             & Schema abuse, dependency-chain injection, batched-auth bypass, IDOR                         \\
		Multi-host / Active Directory & Incalmo MHBench (40 environments) \citep{singer2025incalmo}                                                                                                                                            & Lateral movement, credential reuse, privilege escalation                                    \\
		\bottomrule
	\end{tabular}
\end{table}

\textbf{What is currently left out.} General REST API exploitation is not assessed directly: no standardised, reusable REST API benchmark comparable to CVE-Bench or PrediQL appears in the reviewed literature. Some underlying techniques from the REST API testing literature (dependency inference between API calls, response-feedback pruning) may still be useful internally, but REST-API-specific results are not part of the plan. Binary exploitation, physical- and network-layer attacks, and social engineering are also outside the current scope, for the same reason --- the absence of a suitable existing benchmark in the reviewed literature.
